\chapter{流式 Resume 状态机（跨代对照）}
\label{ch:streaming}

\section{共同契约}

所有后三代流式路径（512\,KiB feed）必须满足：
\begin{equation}
\mathrm{Cuts}^{\mathrm{whole}}(F)=\mathrm{Cuts}^{\mathrm{stream}}(F).
\end{equation}
外层仍由 \texttt{Process*Chunk}：维护当前 chunk 起点、扫描窗、是否允许尾切。

\section{Hom Resume}

\begin{lstlisting}[caption={TopoHomResume},label={lst:hom-resume-full}]
struct TopoHomResume {
  uint32_t h{0};
  size_t scan_rel{0};
  bool in_scan{false};
  bool uf_ready{false};   // production often rebuilds on primary
  SlotUfWindow uf{};
  uint32_t window_w{64};
};
\end{lstlisting}

跨 feed：恢复 $h$ 与相对扫描点；UF 可重建。数学保证：下一字节的 $h$ 递推仅依赖窗尾 $w$ 字节 + 当前 $h$。

\section{Chain Resume}

\begin{lstlisting}[caption={TopoChainResume（含 $S$）},label={lst:chain-resume-full}]
struct TopoChainResume {
  uint32_t S{0};
  size_t scan_rel{0};
  bool in_scan{false};
  bool window_ready{false};
  uint32_t window_w{64};
  uint32_t edge_diff{0};
  uint8_t keys[kChainMaxWindow]{};
};
\end{lstlisting}

与 Hom 差异：可信任并恢复 LFSR $S$（及 keys），并行分段亦依赖串行预热同一递推。跨 feed 若有 carry：\texttt{force\_serial}。

\section{Gen5 TopoPH Resume}

\begin{lstlisting}[caption={TopoPhResume},label={lst:ph-resume}]
struct TopoPhResume {
  uint32_t h{0};
  size_t scan_rel{0};
  bool in_scan{false};
  uint32_t window_w{64};
  uint32_t recent_h[24], recent_t[24];
  uint32_t recent_count{0};
  uint32_t prev_flip{0};
  uint32_t prev_beta0[5]{};
  bool has_prev{false};
};
\end{lstlisting}

Landmark 在 primary 时推入；事件差分依赖 \texttt{has\_prev}。流式须完整搬运 recent 环。

\section{Native Resume 与切后复位}

\begin{align}
&\textbf{跨 feed：} \text{搬运 recent 环、prev\_flip / prev\_beta0、last\_cut}\\
&\textbf{切后：} \texttt{ResetLandmarkWindow}\Rightarrow\text{环清空}
\end{align}

\begin{designbox}[为何切后必须复位]
若不复位，下一 chunk 继承「插入前」的 landmark 几何，1-byte 移位后无法重新锁边，$\mathrm{reuse\_1byte}$ 崩溃。代码注释明确此约束（\filepath{topo\_phn\_scan.cc}）。
\end{designbox}

\section{状态机简图}

\begin{figure}[htbp]
\centering
\begin{tikzpicture}[
  st/.style={rectangle, rounded corners=2pt, draw=blue!50!black, fill=blue!6,
    align=center, font=\small, minimum width=2.4cm, minimum height=0.75cm},
  arr/.style={-{Stealth}, thick}]
  \node[st] (idle) {idle};
  \node[st, right=1.2cm of idle] (scan) {in\_scan};
  \node[st, right=1.2cm of scan] (cut) {emit cut};
  \node[st, below=1.0cm of scan] (feed) {feed boundary};
  \draw[arr] (idle)--node[above,font=\tiny]{enter window}(scan);
  \draw[arr] (scan)--node[above,font=\tiny]{hit}(cut);
  \draw[arr] (cut)--node[right,font=\tiny]{reset?/new pos}(idle);
  \draw[arr] (scan)--(feed);
  \draw[arr] (feed)--node[below,font=\tiny]{restore resume}(scan);
\end{tikzpicture}
\caption{跨代共用的扫描状态机骨架。}
\label{fig:stream-sm}
\end{figure}
